%% ================================================================================================
%%  4. What research says about veto rights
%%  Notes: ../literature/consensus-2026-10_review-notes.md (the Consensus review), research-agenda.md §4
%% ================================================================================================
\section{What research says about veto rights}\label{sec:veto}

The literature on veto rights is split across three traditions that rarely cite each other: committee
and legislative bargaining (economics and political science), the choice of a compromise from a finite
set of options (social choice), and the study of consensus in real organisations. We review each, then
summarise the evidence in Table~\ref{tab:evidence}.

\subsection{Unanimity against majority with a veto}

The sharpest comparison concerns committees that decide between a reform and the status quo.
\citet{BoutonLlorenteSaguerMalherbe2018} compare unanimity with \emph{majority rule with veto power}, in
which members vote by majority but each may also cast a veto that blocks the reform. When members share a common interest but hold private information, the second rule
aggregates information better and dominates unanimity in large committees: under unanimity a member
who doubts the reform can only block it, whereas majority with veto lets them vote against it without
blocking and keeps the veto for the case where they are sure. A laboratory study confirms the ranking
\citep{BoutonLlorenteSaguerMalherbe2017}. The lesson the authors draw, and the Consensus review
endorses, is to separate the right to \emph{block an unacceptable outcome} from the right to
\emph{approve a choice}.

Bargaining models and experiments point the same way. Unanimity can be uniquely inefficient, with
costly conflict and delay \citep{MiettinenVanberg2025,PiazoloVanberg2025}; laboratory committees show
delay and a larger share captured by members who hold a veto
\citep{LaingSlotznick1991,KagelSungWinter2010,AgranovTergiman2019}. The case for unanimity is
conditional: it can improve the honesty of deliberation when members are held responsible for their
votes \citep{BreitmoserValasek2024}, and the comparison with majority depends on what proposers know
\citep{GlyniaThumXefteris2026}.

\subsection{Veto players, agenda power and the status quo}

In political systems, each additional veto player, and each increase in the distance between their
positions, shrinks the set of policies that can replace the status quo; policy stability, or gridlock,
follows \citep{Tsebelis1995,Tsebelis2002}. The status quo is not a neutral fallback: the worse it is for
the others, the more an agenda setter can extract \citep{RomerRosenthal1978}, and whoever proposes holds
power that grows with the patience of the others and the voting rule \citep{BaronFerejohn1989,Harrington1990}.
When today's decision becomes tomorrow's status quo, the dynamics compound: an endogenous status quo
produces inertia and inefficiency \citep{DziudaLoeper2016,AnesiSeidmann2015}, and in laboratory
committees extra rounds strengthen the vetoer \citep{Nunnari2021}; on the stability of policies with
veto players, see \citet{AnesiDuggan2017}. A veto also
discourages experimentation, since a member who may lose from a reform can block the trial that would
reveal its value \citep{AnesiBowen2021}. The best use of a veto player is sometimes to delegate the
choice to them within an interval fixed in advance \citep{KartikKleinerVanWeelden2021}.

Two implications for our design follow. First, the status quo should not be a privileged default:
giving it that role turns every veto into a vote for inertia. Second, the agenda should not be
sequential: in sequential and amendment procedures, the order in which options are put to the vote can be
manipulated \citep{McKelvey1976,BarberaGerber2017}. Under H2, voters judge all options at once, which
removes the agenda setter's power, provided the rule is robust to the composition of the agenda
(Section~\ref{sec:aggregation}).

\subsection{How many rounds?}

There is no general result that more rounds are better. In proposal models with single-peaked
preferences, recent procedures reach the Condorcet winner in one round, where each voter either proposes
or votes \citep{Gersbach2025propose}, or in two rounds with randomly drawn proposers
\citep{GersbachValkanova2024}; see also \citet{Gersbach2024forms}. When voting is costly, a sequence of
rounds with a falling threshold can be optimal \citep{Bognar2015}. In bargaining with vetoes and an
endogenous status quo, however, extra rounds tend to create delay and to strengthen vetoers rather than
improve decisions \citep{Nunnari2021,DziudaLoeper2016}. The Consensus review concludes that one or two
rounds suffice in structured settings, and that robustness of fast procedures to hard vetoes is an open
question. Our design uses one ballot round, with at most one revision round (Section~\ref{sec:safeguards}).

\subsection{Compromise rules with vetoes}

Social choice offers several procedures that select a compromise from a finite set of options by
vetoing.
\begin{itemize}
  \item \textbf{Voting by veto} \citep{Mueller1978}: each of $n$ members proposes one option, the status
  quo is added, and in a random order each member strikes one option; the last one standing wins. The
  procedure always decides and no member gets their worst outcome, but it is not anonymous and needs
  exactly one more option than voters.
  \item \textbf{Proportional veto} \citep{Moulin1981,Moulin1982,Moulin1983}: coalitions veto a share of
  the options equal to their share of the voters (Definition~\ref{def:pvc}). The core is never empty,
  and procedures of successive vetoes can implement it in strong equilibrium. The core can be computed,
  and an option selected from it by an anonymous and neutral rule, in polynomial time
  \citep{IanovskiKondratev2021}; \citet{KondratevIanovski2024} study which rules that always select from
  the core are monotonic, satisfy participation or ignore options that everyone ranks last.
  \item \textbf{Fallback bargaining} \citep{BramsKilgour2001}: members descend their rankings in
  lockstep until an option appears on every list. The outcome is Pareto-optimal and maximises the worst
  rank, but it may miss the Condorcet winner and even a majority's first choice; with an ``impasse'' line,
  below which members refuse to go, it can end in no agreement.
  \item \textbf{Majoritarian compromise} \citep{SertelYilmaz1999,HurwiczSertel1999}: the same descent,
  stopped at the first depth at which an option appears in the lists of a majority; it is
  majoritarian-optimal and implementable in subgame-perfect equilibrium.
  \item \textbf{Shortlisting and veto-rank} for two parties: one party shortlists, the other chooses
  \citep{BarberaCoelho2017,BarberaCoelho2022}; or parties strike options alternately
  \citep{Anbarci1993}, or in proportion to ranks \citep{LaslierNunezSanver2021}. These rules can have
  several equilibria, some Pareto-inefficient, and \citet{CaillouxNapolitanoSanver2022} argue that several
  ``compromise'' rules fail to select an outcome that is a compromise once the result is known.
\end{itemize}
None of these was designed for an open agenda. Section~\ref{sec:aggregation} shows that fallback
bargaining, the majoritarian compromise, sequential vetoes and the proportional veto core all reward
flooding the agenda with clones.

\subsection{Consensus in practice}

Organisations that decide ``by consensus'' rarely give every member a hard veto. The Internet
Engineering Task Force works by \emph{rough consensus}: lack of disagreement matters more than agreement,
and consensus is reached when all objections have been addressed, though not necessarily accommodated;
an objection judged unpersuasive does not block \citep[\S\S2--3]{Resnick2014rfc7282}. In bodies that
decide by tacit consent, the formal right of every member to object is constrained by social pressure
and by the prospect of a fallback majority vote, so that effective veto power is much smaller than the
rule suggests \citep{Krick2017}. Where a hard individual veto has been taken literally, as in the
\emph{liberum veto} of the Polish--Lithuanian Sejm, it paralysed decision-making \citep{Rohac2008}.
Formal institutions grant vetoes to coalitions instead: the Council of the European Union's blocking
minority needs at least four members \citep{TEU2016art16}.

The closest existing procedure to the one we recommend is the Debian constitution
\citep[\S A.5]{Debian2024constitution}. Developers rank all the options of a general resolution,
including the default option ``None of the above''; options ranked above the default are deemed
acceptable. An option is dropped if too few voters rank it above the default (a quorum) or if it fails
to beat the default by its required majority (3:1 for constitutional amendments); the winner among the
remaining options is chosen by the Schulze method. Debian thus combines open amendments, full ranked
ballots, an acceptability line and a clone-independent Condorcet method, but it keeps the default
option as a privileged fallback, a choice that suits a project in which ``no change'' is a safe
outcome.

\begin{table}[htbp]\centering\small
\caption{What the literature supports}\label{tab:evidence}
\begin{tabularx}{\linewidth}{>{\raggedright\arraybackslash}X l >{\raggedright\arraybackslash}p{0.33\linewidth}}
\toprule
Claim & Evidence & Main sources \\
\midrule
Full unanimity is rarely optimal for open choices: it raises delay, status-quo bias and holdout power
  & Strong & \citet{Tsebelis2002,DziudaLoeper2016,MiettinenVanberg2025,KagelSungWinter2010}; Prop.~\ref{prop:gridlock} \\
Majority with a separate veto beats unanimity for information aggregation
  & Strong & \citet{BoutonLlorenteSaguerMalherbe2018,BoutonLlorenteSaguerMalherbe2017} \\
No rule is strategy-proof; the majority principle and minority vetoes conflict
  & Theorem & \citet{Gibbard1973,Satterthwaite1975}; Prop.~\ref{prop:majority-veto} \\
Simultaneous voting on all options removes agenda-setter power; clone independence removes the gain
  from flooding & Strong (theory) & \citet{McKelvey1976,Tideman1987,Schulze2011} \\
Condorcet winners exist in most structured electorates
  & Strong & \citet{RegenwetterEtAl2006,TidemanPlassmann2012}; Table~\ref{tab:cw_probability} \\
Fallback, shortlisting and veto-rank rules are useful second-best compromise devices
  & Moderate & \citet{BramsKilgour2001,BarberaCoelho2022,CaillouxNapolitanoSanver2022} \\
One or two rounds suffice in structured proposal settings
  & Moderate & \citet{GersbachValkanova2024,Gersbach2025propose} \\
In practice, a formal individual veto is socially constrained
  & Weak to moderate & \citet{Krick2017,Resnick2014rfc7282} \\
\bottomrule
\end{tabularx}
\figurenotes{Evidence: strong, several formal results and experiments or field evidence agree;
moderate, formal results under restrictive assumptions; weak, case studies. Adapted and extended from
the Consensus review of October 2026 (its Figure~6).}{Sections~\ref{sec:limits}--\ref{sec:aggregation}
of this note.}
\end{table}
